% !TEX root = task-2-10-5.tex
% Editable task bank. Append new BankTask blocks at the end of the document.
\documentclass[12pt,a4paper]{article}
\usepackage{physics-textbook}
\usepackage{graphicx,needspace}
\usetikzlibrary{decorations.pathmorphing}
\newif\ifShowSolutions
\ifdefined\TutorEdition
  \ShowSolutionstrue
\else
  \ShowSolutionsfalse % Change to \ShowSolutionstrue to display solutions.
\fi
\newcounter{banktask}
\newsavebox{\BankTaskBox}
\newcommand{\BankTaskContent}[3]{%
  \noindent\textbf{\thebanktask.}\quad #2\par
  \ifShowSolutions
    \Needspace{50mm}% Keep the solution heading with its opening text or diagram.
    #3\par
    \ifdefined\TutorEdition
      \ifcsname TutorSource#1\endcsname
        \nopagebreak\smallskip{\small\csname TutorSource#1\endcsname\par}%
      \fi
    \fi
  \fi
}
\newcommand{\BankTask}[3]{%
  \refstepcounter{banktask}\par\medskip
  \sbox{\BankTaskBox}{\begin{minipage}{\linewidth}%
    \BankTaskContent{#1}{#2}{#3}\end{minipage}}%
  \ifdim\dimexpr\ht\BankTaskBox+\dp\BankTaskBox\relax>\dimexpr\textheight-20mm\relax
    \Needspace{35mm}\BankTaskContent{#1}{#2}{#3}%
  \else
    \Needspace{\dimexpr\ht\BankTaskBox+\dp\BankTaskBox\relax}%
    \noindent\usebox{\BankTaskBox}\par
  \fi
  \medskip
}
\begin{document}
\BookSubsection{2.10.5}{Парообразование и конденсация}


% Source sdamgia-30492: https://phys-ege.sdamgia.ru/problem?id=30492
\BankTask{sdamgia-30492}{
Какое количество теплоты нужно сообщить при нормальном атмосферном давлении воде массой $0{,}5\,\text{кг}$, уже нагретой до температуры $100^\circ$С, для того, чтобы она полностью выкипела? (В килоджоулях.)
}{
\smallskip\noindent\textbf{Решение}\par
При кипении требуется количество теплоты, равное:

$Q=Lm=2{,}3 \cdot 10 ^{6} \cdot 0{,}5=1150\text{кДж}.$

Ответ: $1150$.
}

% Source sdamgia-31882: https://phys-ege.sdamgia.ru/problem?id=31882
\BankTask{sdamgia-31882}{
Какое количество теплоты нужно сообщить $2\,\text{кг}$ воды, нагретой до температуры $50^\circ$C, чтобы она полностью выкипела? (В килоджоулях.)
}{
\smallskip\noindent\textbf{Решение}\par
Для нагревания воды от начальной температуры до температуры кипения потребуется количество теплоты:

$Q_1=cm ( t_{\text{к}} - t ) .$

При кипении потраченное количество теплоты равно $Q_2=Lm.$ Тогда общее количество теплоты равно:

$Q=Q_1 + Q_2=m ( c ( t_{\text{к}} - t ) + L ) =2 \cdot ( 4200 \cdot 50 + 2{,}3 \cdot 10 ^{6} ) =5{,}02 \cdot 10 ^{6} \text{Дж}=5020\text{кДж}.$

Ответ: $5020$.

Примечание. 

Значения удельной теплоемкости и удельной теплоты парообразования берем из справочника к КИМам на нашем сайте (правый верхний угол — круг).
}

% Source sdamgia-43491: https://phys-ege.sdamgia.ru/problem?id=43491
\BankTask{sdamgia-43491}{
Воде, нагретой до температуры $100 ^\circ C$ при нормальном атмосферном давлении, сообщили количество теплоты $Q = 4{,}6 \cdot 10 ^{6} \text{Дж}.$ Сколько воды при этом испарится? (В килограммах.)
}{
\smallskip\noindent\textbf{Решение}\par
Так как вода уже нагрета до температуры кипения, то при процессе парообразования $Q=Lm,$ откуда масса испарившейся воды равна

$m= \frac{Q}{L} = \frac{4{,}6 \cdot 10 ^{6}}{2{,}3 \cdot 10 ^{6}} =2\text{кг}.$

Ответ: $2$.
}

% Source sdamgia-7348: https://phys-ege.sdamgia.ru/problem?id=7348
\BankTask{sdamgia-7348}{
Утром после восхода солнца роса, лежавшая на траве, через некоторое время исчезла. Это произошло потому, что

$1$) водяной пар, который содержался в приземном слое воздуха до восхода солнца, был ненасыщенным.

$2$) после восхода солнца температура в приземном слое воздуха повысилась, и водяной пар, который содержался в воздухе, стал ненасыщенным.

$3$) относительная влажность приземного слоя воздуха после восхода солнца увеличилась.

$4$) после восхода солнца трава начала менее интенсивно выделять содержащуюся в ней воду, и капли воды с травинок испарились.
}{
\smallskip\noindent\textbf{Решение}\par
Это произошло потому, что после восхода солнца температура в приземном слое воздуха повысилась, и водяной пар, который содержался в воздухе, стал ненасыщенным.

Ответ: $2$
}

% Source sdamgia-3372: https://phys-ege.sdamgia.ru/problem?id=3372
\BankTask{sdamgia-3372}{
На газовой плите стоит высокая кастрюля с водой, закрытая крышкой. Если воду из нее перелить в широкую кастрюлю, у которой площадь дна вдвое больше, и тоже закрыть крышкой, то вода закипит заметно быстрее, чем если бы она осталась в узкой. Этот факт объясняется тем, что

$1$) увеличивается площадь нагревания и, следовательно, увеличивается скорость нагревания воды

$2$) в $2$ раза уменьшается необходимое давление насыщенного пара в пузырьках и, следовательно, воде у дна надо нагреваться до менее высокой температуры

$3$) увеличивается площадь поверхности воды и, следовательно, испарение идет более активно

$4$) в $2$ раза уменьшается глубина слоя воды и, следовательно, пузырьки пара быстрее добираются до поверхности
}{
\smallskip\noindent\textbf{Решение}\par
Процесс кипения начинается при температуре, когда внешнее давление на жидкость сравнивается с давлением насыщенных паров жидкости при этой температуре. От площади поверхности жидкости температура кипения никак не зависит То, что вода закипит в более широкой кастрюле быстрее, объясняется тем, что при этом увеличивается площадь нагревания, и, следовательно, увеличивается скорость нагревания воды.

Ответ: $1$.
}

% Source sdamgia-2922: https://phys-ege.sdamgia.ru/problem?id=2922
\BankTask{sdamgia-2922}{
В цилиндрическом сосуде под поршнем длительное время находятся вода и ее пар. Поршень начинают медленно выдвигать из сосуда. При этом температура воды и пара остается неизменной. Как будет меняться при этом масса жидкости в сосуде? Ответ поясните, указав, какие физические закономерности вы использовали для объяснения.
}{
\smallskip\noindent\textbf{Решение}\par
$1$. Ответ: масса жидкости в сосуде будет уменьшаться.

$2$. Вода и водяной пар находятся в закрытом сосуде длительное время, поэтому водяной пар является насыщенным.

$3$. При выдвигании поршня пар изотермически расширяется. Давление и плотность насыщенного пара в этом процессе не меняются. Следовательно, для пополнения количества вещества пара будет происходить испарение жидкости. Значит, масса жидкости в сосуде будет уменьшаться.
}

% Source sdamgia-3674: https://phys-ege.sdamgia.ru/problem?id=3674
\BankTask{sdamgia-3674}{
В герметичную банку, сделанную из очень тонкой жести и снабженную наверху завинчивающейся крышкой, налили немного воды (заполнив малую часть банки) при комнатной температуре и поставили на газовую плиту, на огонь, не закрывая крышку. Через некоторое время, когда почти вся вода выкипела, банку сняли с огня, сразу же плотно завинтили крышку и облили банку холодной водой. Опишите физические явления, которые происходили на различных этапах этого опыта, а также предскажите и объясните его результат.
}{
\smallskip\noindent\textbf{Решение}\par
$1$. После помещения банки на огонь вода в ней нагревалась через тонкую стенку банки от горячих продуктов горения газа. При этом с ростом температуры вода испарялась, и возрастало давление ее паров в банке, которые постепенно вытесняли из нее воздух. Когда вода закипела и почти вся испарилась, воздуха внутри банки практически не осталось. Давление насыщенных паров в банке при этом стало равно внешнему атмосферному давлению.

$2$. Когда банку сняли с огня, закрыли крышкой и охладили холодной водой почти до комнатной температуры, горячие пары воды внутри банки остыли и практически целиком сконденсировались на ее стенках, отдавая теплоту конденсации наружу, холодной воде, благодаря процессу теплопроводности через стенки.

$3$. В соответствии с уравнением Клапейрона — Менделеева давление пара в банке резко упало — во-первых, из-за уменьшения массы оставшегося в банке пара, и во-вторых — из-за падения его температуры. Заметим, что резкое уменьшение давления в банке можно объяснить и так: при понижении температуры до комнатной пары конденсируются, оставаясь насыщенными, но их давление становится намного меньше давления насыщенных паров воды при температуре кипения (примерно в $40$ раз).

$4$. Поскольку при комнатной температуре давление насыщенных паров воды составляет лишь малую долю от атмосферного давления (не более $3$-$4$\%), тонкая банка после поливания ее водой окажется под действием разности этого большого внешнего давления и низкого давления пара внутри. По этой причине на банку начнут действовать большие сдавливающие силы, которые будут стремиться сплющить банку. Как только эти силы превысят предельную величину, которую могут выдержать стенки банки, то она сплющится, резко уменьшившись в объеме.
}

% Source sdamgia-38577: https://phys-ege.sdamgia.ru/problem?id=38577
\BankTask{sdamgia-38577}{
В кастрюле находится $0{,}5\,\text{кг}$ воды при температуре $10 ^\circ\text{C}$. Сколько потребуется времени, чтобы при помощи кипятильника с постоянной потребляемой мощностью $400\,\text{Вт}$ превратить в пар $15\%$ воды из кастрюли? Потерями тепла и теплоёмкостью кастрюли пренебречь.
}{
\smallskip\noindent\textbf{Решение}\par
Вода в кастрюле нагревается до температуры кипения $100^\circ\text{C}$, получив количество теплоты
\[Q_1=mc(t_{\text{к}}-t),\]
а затем частично превращается в пар, получив количество теплоты
\[Q_2=0{,}15mL.\]
Так как потерь тепла не происходило, то это количество теплоты равно количеству теплоты, полученному от кипятильника. Из формулы тепловой мощности $P=\dfrac{Q}{\tau}$ при объединении формул получаем:
\[P=\frac{m\bigl(c(t_{\text{к}}-t)+0{,}15L\bigr)}{\tau},\]
откуда находим необходимое время:
\[\begin{aligned}
\tau&=\frac{m\bigl(c(t_{\text{к}}-t)+0{,}15L\bigr)}{P}\\
&=\frac{0{,}5\bigl(4{,}2\cdot10^3(100-10)+0{,}15\cdot2{,}3\cdot10^6\bigr)}{400}\approx904\,\text{с}.
\end{aligned}\]
Ответ: $904\,\text{с}$.
}

% Source sdamgia-25410: https://phys-ege.sdamgia.ru/problem?id=25410
\BankTask{sdamgia-25410}{
В вертикальном цилиндре, закрытом легким поршнем, находится бензол ($\mathrm{C_6H_6}$) при температуре кипения t  =  $80$ \ensuremath{{}^\circ}C. При сообщении бензолу количества теплоты Q часть его превращается в пар, который при изобарном расширении совершает работу А. Удельная теплота парообразования бензола L  =  $396$ · $10^{3}$ Дж/кг, его молярная масса M  =  $78$ · $10^{-3}$ кг/моль. Какая часть подведенного к бензолу количества теплоты переходит в работу? Объемом жидкого бензола пренебречь.
}{
\smallskip\noindent\textbf{Решение}\par
$1$.   При кипении бензола происходит его изобарное расширение. Работа пара $A = p\Delta V,$ где p  — атмосферное давление, $\Delta V$  — изменение объема.
$2$.  Считая пар идеальным газом, воспользуемся уравнением Клапейрона  — Менделеева для определения изменения объема за счет испарившегося бензола массой $\Delta m:$ $p\Delta V= \frac{\Delta m}{M} RT,$ где M  =  $78$ · $10^{-3}$ кг/моль  — молярная масса бензола, T  =  $80$ + $273$  =  $353$ K  — температура кипения бензола. Отсюда $A= \frac{\Delta mRT}{M} .$
$3$.  Количество теплоты Q, необходимое для испарения массы $\Delta m$ бензола, пропорционально удельной теплоте парообразования L: $Q = \Delta mL .$
$4$.  Искомая величина определяется отношением:

\[\eta = \frac{A}{Q} = \frac{RT}{ML} = \frac{8{,}31 \cdot 353}{78 \cdot 10 ^{ - 3 } \cdot 396 \cdot 10 ^3} \approx 0{,}095.\]

Ответ:$\eta \approx 0{,}095.$
}

% Source sdamgia-42650: https://phys-ege.sdamgia.ru/problem?id=42650
\BankTask{sdamgia-42650}{
\begin{center}\begin{tikzpicture}[task graph,>=stealth]\draw[gray!65,xstep=.28,ystep=.29](0,0)grid(5.6,4.64);
\draw[->](0,0)--(6.8,0)node[below]{\(t,\,^\circ\text{C}\)};
\draw[->](0,0)--(0,5.2)node[left]{\(p_{\text{н}},\,\text{гПа}\)};
\node[below left]at(0,0){\(0\)};
\foreach \x/\label in {1.4/10,2.8/20,4.2/30,5.6/40}{\node[below]at(\x,0){\(\label\)};}
\foreach \y/\label in {1.16/20,2.32/40,3.48/60,4.64/80}{\node[left]at(0,\y){\(\label\)};}
\draw[smooth] plot coordinates{(0,.47)(.7,.66)(1.4,.83)(2.24,1.16)(2.8,1.53)(3.5,2.1)(4.2,2.78)(4.9,3.65)(5.53,4.99)};\end{tikzpicture}\end{center}
В сосуде ёмкостью V  =  $2$ л, воздух из которого был откачан до высокого вакуума, разбили заполненную водой ампулу ёмкостью $V_{1}$  =  $4\,\text{см}^{3}$. Какая часть воды, содержавшейся в ампуле, испарилась? В сосуде поддерживалась температура $16 ^\circ\text{C}.$ Зависимость давления насыщенного пара от температуры представлена на графике. Объёмом воды в сосуде можно пренебречь.
}{
\smallskip\noindent\textbf{Решение}\par
После того, как в сосуде разбили пробирку, в нем оказалась вода массой $m_{\,\text{в}}=\rho V_{\,\text{п}}=10 ^3 \cdot 4 \cdot 10 ^{ - 6 } =4\,\text{г}.$ Поскольку в сосуде находится высокий вакуум, то при температуре $16$ \ensuremath{{}^\circ}C вода закипит и превратится в пар. Если пар в сосуде стал насыщенным, то, согласно графику, его давление должно быть равным $20$ гПа.
При температуре $16$ \ensuremath{{}^\circ}C ($289$ К) сосуде в результате испарения воды оказался пар, массу  которого найдем из уравнения Менделеева-Клапейрона $pV= \frac{m}{M} RT$ при условии, если бы пар стал насыщенным:

\[m= \frac{pVM}{RT} = \frac{2 \cdot 10 ^3 \cdot 2 \cdot 10 ^{ - 3 } \cdot 18 \cdot 10 ^{ - 3 }}{8{,}31 \cdot 289} \approx 30 \cdot 10 ^{ - 3 } \,\text{г}.\]

Отсюда делаем вывод, что пар стал насыщенным, а из пробирки испарилась часть воды, равная

\[\frac{m}{m_{\,\text{в}}} = \frac{30 \cdot 10 ^{ - 3 }}{4} =0{,}0075.\]

Ответ: $0{,}0075$.
}


% Source 6209: https://phys-ege.sdamgia.ru/problem?id=6209
\BankTask{6209}{
В закрытом сосуде находится $2\,\text{г}$ водяного пара под давлением $50$ кПа и при температуре $100$ ºС. Не изменяя температуры, объем сосуда уменьшили в $4$ раза. Найдите массу образовавшейся при этом воды. Ответ приведите в граммах.
}{
\smallskip\noindent\textbf{Решение}\par
При уменьшении объема сосуда будет увеличиваться давление газа. Достигнув давления насыщенного пара давление перестанет меняться. При дальнейшем уменьшении объема сосуда давление перестанет меняться, но из газа начнет конденсироваться жидкость. Вспомним, что вода кипит при температуре $100$ \ensuremath{{}^\circ}C и нормальном давлении $10^{5}$ Па. Следовательно, давление насыщенного пара при температуре $100$ \ensuremath{{}^\circ}C для воды составляет $100$ кПа. Напишем уравнение Клапейрона  — Менделеева для первого и второго состояний газа.

\[p_1V_1= \frac{m_1}{\mu} RT,p_2V_2= \frac{m_2}{\mu} RT,\]

где $m_1$  — масса газа в первом состоянии, $m_2$  — масса газа во втором состоянии. Из этих соотношений находим массу газа во втором состоянии:

\[m_2= \frac{p_2V_2}{p_1V_1} m_1= \frac{p_2V_2}{p_1 \cdot 4V_2} = \frac{100}{50 \cdot 4} \cdot 2=1\,\text{г}.\]

Масса пара во втором состоянии равна $1\,\text{г}$, следовательно, масса образовавшейся воды $2$ − $1$  =  $1\,\text{г}$.

Ответ: $1\,\text{г}$.
}

% Source 6281: https://phys-ege.sdamgia.ru/problem?id=6281
\BankTask{6281}{
В сосуде под поршнем при температуре $100$ ºС находится $2\,\text{г}$ водяного пара и такое же количество воды. Не изменяя температуры, объем сосуда увеличили в $3$ раза. Определите массу воды, перешедшей при этом в пар. Ответ приведите в граммах.
}{
\smallskip\noindent\textbf{Решение}\par
Под поршнем находится насыщенный пар, то есть жидкость и газ находятся в динамическом равновесии друг с другом. При каждом небольшом увеличении объема сосуда давление газа немного уменьшается. В результате этого часть молекул воды переходит из жидкого состояния в газообразное, динамическое равновесие восстанавливается. Поэтому, пока вся вода не перейдет в газообразное состояние, давление пара над жидкостью не будет меняться. Вспомним, что вода кипит при температуре $100$ °C и нормальном давлении $10^{5}$ Па. Следовательно, давление насыщенного пара при температуре $100$ °C для воды составляет $100$ кПа. Напишем уравнение Клапейрона  — Менделеева для первого и второго состояний газа.

\[p_1V_1= \frac{m_1}{\mu} RT,p_2V_2= \frac{m_2}{\mu} RT,\]

где $m_1$  — масса газа в первом состоянии, $m_2$  — масса газа во втором состоянии. Из этих соотношений находим массу газа во втором состоянии:

\[m_2= \frac{p_2V_2}{p_1V_1} m_1= \frac{p_1 \cdot 3V_1}{p_1V_1} m_1= \frac{3}{1} \cdot 2=6\,\text{г}.\]

В результате расчетов получается, что в газообразное состояние должны были бы перейти $6\,\text{г} - 2\,\text{г}=4\,\text{г}$ воды. Учитывая, что воды в жидком состоянии лишь $2\,\text{г}$, получаем что вся вода перейдет в газообразное состояние, то есть масса воды, перешедшей в пар, равна $2\,\text{г}$.

Ответ: $2\,\text{г}$.
}

% Source 3721: https://phys-ege.sdamgia.ru/problem?id=3721
\BankTask{3721}{
В закрытом сосуде под поршнем находится $4\,\text{г}$ насыщенного водяного пара. Двигая поршень, занимаемый паром объем уменьшили в $2$ раза, поддерживая температуру сосуда и его содержимого постоянной и равной $100$ °C. Какое количество теплоты было при этом отведено от сосуда?
}{
\smallskip\noindent\textbf{Решение}\par
Насыщенный пар  — это пар, имеющий максимально возможную концентрацию при данной температуре, при этом концентрация пара зависит только от температуры. Поскольку при уменьшении объема в два раза, температуру сосуда и его содержимого поддерживали постоянной, заключаем, что концентрация насыщенного пара не изменялась, а значит, в результате такого сжатия сконденсировалась ровно половина всего пара, то есть $2\,\text{г}$.
По первому началу термодинамики изменение внутренней энергии системы равно разности работы, совершаемой над системой, и количеством теплоты, забираемой у системы:

$\Delta U=A - Q ,$ откуда $Q= - \Delta U + A.$

При конденсировании $2\,\text{г}$ пара его внутренняя энергия уменьшается:

\[\Delta U= - L\Delta m= - 2{,}3 \cdot 10 ^{6} \,\text{Дж/кг} \cdot 0{,}002 \,\text{кг}= - 4600\,\text{Дж}.\]

При сжатии над паром совершается работа:

\[A=p\Delta V= \frac{pV_0}{2} = \frac{m}{2M} RT= \frac{4\,\text{г}}{2 \cdot 18\,\text{г/моль}} \cdot 8{,}31\,\text{Дж}/ ( \,\text{моль} \cdot \,\text{К} ) \cdot 373\,\text{К}\approx344\,\text{Дж}.\]

Количество забираемой теплоты равно:

\[Q=4600\,\text{Дж} + 344\,\text{Дж}\approx5\,\text{кДж}.\]

Ответ: $5\,\text{кДж}$.
}

\end{document}
